% Generated draft: requires mathematical and source review.
\documentclass[11pt]{article}
\usepackage[margin=25mm]{geometry}
\usepackage{fontspec}
\setmainfont{Latin Modern Roman}
\usepackage{amsmath,amssymb,mathtools,unicode-math}
\setmathfont{Latin Modern Math}
\usepackage{enumitem,needspace,xurl,hyperref}
\setlength{\parindent}{0pt}
\setlength{\parskip}{5pt}
\setlength{\emergencystretch}{3em}
\newcommand{\N}{\mathbb{N}}
\newcommand{\Z}{\mathbb{Z}}
\newcommand{\Q}{\mathbb{Q}}
\newcommand{\R}{\mathbb{R}}
\newcommand{\C}{\mathbb{C}}
\newcommand{\F}{\mathbb{F}}
\newcommand{\A}{\mathbb{A}}
\newcommand{\PP}{\mathbb{P}}
\newcommand{\E}{\mathbb{E}}
\newcommand{\eps}{\varepsilon}
\newcommand{\dd}{\,\mathrm{d}}
\begin{document}
\section*{UnsolvedMath ID: 800779}
\textbf{Source problem code:} catalogue entry
\subsection{Definitions and mathematical statement}
The supplied excerpt concerns a holomorphic local dynamical system at a point \(p\), but does not explicitly define either that term or \(K_f\). The following formulation is conditional on the usual interpretation: \(M\) is a finite-dimensional complex manifold, \(U\subseteq M\) is an open neighbourhood of \(p\), and \(f\colon U\to M\) is holomorphic with \(f(p)=p\). Define the forward-staying set by
\[
K_f(U)=\{x\in U: f^j(x)\text{ is defined and belongs to }U\text{ for every integer }j\geq 0\}.
\]
Here \(f^0(x)=x\), and subsequent iterates are defined recursively while their inputs lie in \(U\). Equip \(K_f(U)\) with the topology induced by \(M\). This interpretation makes \(f\colon K_f(U)\to K_f(U)\) a dynamical system. The excerpt does not confirm this definition of \(K_f\), fix the neighbourhood \(U\), or specify how the requested description should depend on shrinking \(U\).

\textbf{Problem Q2, conditional on this interpretation.} For a given holomorphic local dynamical system \(f\) at \(p\), determine the dynamical structure of \((K_f(U),f)\). In particular, describe the asymptotic behaviour of the forward orbits \((f^j(x))_{j\geq 0}\), for \(x\in K_f(U)\). Determine which such orbits converge to \(p\), and whether other behaviour, described in the source as chaotic, occurs. Determine whether there exists \(x\in K_f(U)\) such that
\[
\overline{\{f^j(x):j\geq 0\}}^{\,K_f(U)}=K_f(U).
\]
Describe proper forward-invariant subsets \(L\subsetneq K_f(U)\), meaning subsets satisfying \(f(L)\subseteq L\), and the restricted dynamics \(f|_L\colon L\to L\). The source does not define chaotic behaviour or impose nonemptiness, closedness, or another nontriviality condition on \(L\); these qualifications remain unresolved.

The excerpt proposes comparison with a simpler system, for example a linear system, by local conjugacy. Its supplied definition is as follows. Let \(f_i\colon U_i\to M_i\) be holomorphic local dynamical systems at \(p_i\), for \(i\in\{1,2\}\). They are holomorphically, smoothly, \(C^k\), or topologically locally conjugated if there are open neighbourhoods \(W_i\subseteq U_i\) of \(p_i\) and a map \(\varphi\colon W_1\to W_2\), respectively a biholomorphism, a \(C^\infty\) diffeomorphism, a \(C^k\) diffeomorphism for a positive integer \(k\), or a homeomorphism, such that \(\varphi(p_1)=p_2\) and
\[
\varphi^{-1}\bigl(W_2\cap f_2^{-1}(W_2)\bigr)
 =W_1\cap f_1^{-1}(W_1),
\qquad
f_1=\varphi^{-1}\circ f_2\circ\varphi
\]
on that common domain. The source states that such a conjugacy identifies the corresponding forward-staying sets:
\[
K_{f_2|_{W_2}}=\varphi\bigl(K_{f_1|_{W_1}}\bigr).
\]
It also explains that local coordinates centred at \(p\) reduce the setting to holomorphic local systems fixing \(0\in\C^n\), where \(n=\dim_{\C}M\), denoted there by \(\operatorname{End}(\C^n,0)\). These are contextual tools for the dynamical-description question; the excerpt does not assert that every system admits a linear model.
\subsection{Short English statement}
Abate's Problem Q2 asks for a description of local holomorphic dynamics: orbit behaviour, dense orbits, and dynamics on proper forward-invariant subsets. The excerpt discusses local conjugacy as a method of comparison. Its notation \(K_f\) and its use of chaotic behaviour are not defined in the supplied record, so the displayed formulation explicitly uses a conditional forward-staying-set interpretation and requires human review.
\subsection{Sources}
\begin{itemize}
\item Local UnsolvedMath registry, supplied complete source record: permanent ID \(800779\), exact source problem code \texttt{catalogue entry}, title \emph{Local Discrete Holomorphic Dynamics --- Open Problem Q2}; collection: AMR Open Problem Lists. The record attributes the problem to Abate, with proposed year \(2011\), and supplies the status label \texttt{open}. This label is source metadata, not an independently verified assessment of the current mathematical status.
\item Abate, \emph{Open Problems in Local Discrete Holomorphic Dynamics} (2011), Problem Q2, supplied source URL: \url{https://arxiv.org/abs/1107.4876}. The record reports access on \texttt{2026-10-16} and extraction method \texttt{arxiv\_latex}. These are supplied provenance claims; the reference and extraction have not been independently checked for this draft.
\end{itemize}
\end{document}
